% Lang Liu -- Resume
% Compile with: pdflatex resume.tex   (or latexmk -pdf resume.tex)
\documentclass[11pt,letterpaper]{article}

\usepackage[T1]{fontenc}
\usepackage[utf8]{inputenc}
\usepackage{newtxtext}            % Times-like font with real small caps
\usepackage[margin=0.5in,top=0.45in,bottom=0.45in]{geometry}
\usepackage{enumitem}
\usepackage{titlesec}
\usepackage{pifont}               % \ding{220} arrow
\usepackage{amssymb}              % \blacksquare
\usepackage{graphicx}
\usepackage[hidelinks]{hyperref}

\pagestyle{empty}
\setlength{\parindent}{0pt}
\setlength{\parskip}{0pt}
\linespread{0.97}

% ---------- Helpers ----------
\newcommand{\sqbullet}{\raisebox{0.3ex}{\scalebox{0.45}{$\blacksquare$}}}
\newcommand{\sep}{\enspace\sqbullet\enspace}
\newcommand{\arrow}{\ding{220}}

% Section heading: bold small caps with a full-width rule underneath
\titleformat{\section}{\fontsize{13.5}{14}\selectfont\bfseries\scshape}{}{0pt}{}[\vspace{-1.6ex}\rule{\textwidth}{0.6pt}]
\titlespacing*{\section}{0pt}{2.2ex}{0.2ex}

% Bullet list
\setlist[itemize]{label=\sqbullet, leftmargin=0.3in, labelsep=0.45em,
  itemsep=0pt, parsep=0pt, topsep=0pt, partopsep=0pt}

% Education entry: {School, Location}{Degree}{Dates}
\newcommand{\edu}[3]{%
  #1\par
  \begin{itemize}\item #2 \hfill #3\end{itemize}}

% Job heading: {Company}{Role}{Dates}
\newcommand{\job}[3]{%
  \textbf{\textsc{#1}} | #2\hfill #3\par}
\newcommand{\jobnote}[1]{\textit{#1}\par}

% Publication: {Title}{Authors and venue}
\newcommand{\pub}[2]{\textit{\textsc{#1}}\par #2\par}

% Make the text wrap justified within bullets like the original
\AtBeginDocument{\sloppy}

\begin{document}

% ---------- Header ----------
{\fontsize{30}{32}\selectfont\bfseries\scshape Lang Liu}%
\hspace{0.6em}\raisebox{0.15ex}{\small\sqbullet\enspace
\href{mailto:langinteger@outlook.com}{langinteger@outlook.com}\sep
\href{https://www.linkedin.com/in/langinteger}{https://www.linkedin.com/in/langinteger}}

% ---------- Education ----------
\section{Education}
\edu{\textsc{Illinois Institute of Technology}, Chicago, IL, United States}
    {Ph.D.\ in Computer Science, GPA: 3.89/4.00}{Jan 2024 - May 2027 (Expected)}
\edu{\textsc{Central South University}, Changsha, Hunan, China}
    {B.S.\ in Civil Engineering, GPA: 86.9/100}{Sep 2013 - Jun 2017}

% ---------- Research Interests ----------
\section{Research Interests}
\textsc{Programming Languages, Formal Methods, Separation Logic, Program Verification, Compiler Security}

% ---------- Publications ----------
\section{Publications}
\pub{Quantified Underapproximation via Labeled Bunches}
    {\textbf{L. \textsc{Liu}}, F. \textsc{Derakhshan}, L. \textsc{Jia}, G. \textsc{Moreno}, M. \textsc{Klein}. OOPSLA 2025}
\pub{End-to-End Guarantees for Compartmentalized Safety Critical Software}
    {\textbf{L. \textsc{Liu}}, Z. \textsc{Zhang}, A. \textsc{Vasudevan}, F. \textsc{Derakhshan}, L. \textsc{Jia}, W. \textsc{Mansky}. \textsc{Under Review}}
\pub{PIMSA: A Static Analysis Framework for Processing-in-Memory Programs}
    {\textbf{L. \textsc{Liu}}, A. \textsc{Madhu}, S. \textsc{Nawaz}, F. \textsc{Derakhshan}, M. \textsc{Zhou}. \textsc{Under Review}}

% ---------- Work Experiences ----------
\section{Work Experiences}
\job{Illinois Institute of Technology}{\textsc{Research Assistant \& Teaching Assistant}}{Jan 2024 -- Present}
\begin{itemize}
  \item Developed logical frameworks for compiler validation and the integration of formal verification and testing.
  \item Teaching Assistant for CS351 System Programming and CS534 Types and Programming Languages.
\end{itemize}

\job{Amazon Web Services}{\textsc{Applied Scientist Intern}}{May 2026 -- Aug 2026}
\jobnote{Automated Reasoning Group}
\begin{itemize}
  \item Developed CranDSL, a Domain-Specific Language (DSL) in Lean 4 for writing test cases for the AWS IAM model.
  \item Leveraged Lean's dependent type system to enforce operation dependencies statically within test cases.
  \item Awarded Top 3 in AWS Identity Hackathon (Operational Excellence category) for building an LLM-aided test
        generation framework using CranDSL, achieving 100\% line coverage for APIs considered.
\end{itemize}

\job{Presence Inc}{\textsc{Founding} Software Engineer}{Apr 2022 -- Aug 2023}
\jobnote{A startup invested by BlueRun Ventures}
\begin{itemize}
  \item Developed backend services for several social apps from inception to launch, supporting over 10k DAU.
  \item Led the agile development process and CI/CD pipeline, ensuring weekly deployments to AWS/GCP and regular app
        releases on Google Play and Apple App Store, improving release efficiency by over 2x.
  \item Tech Stack: Java, Spring Boot, Vert.x, C\#, Unity, Flutter, Redis, PostgreSQL, gRPC, AWS, GCP
\end{itemize}

\job{Tencent Technology (Shenzhen) Co.\ Ltd}{Senior Software Engineer}{Jul 2021 -- Apr 2022}
\jobnote{Cloud \& Smart Industries Business Group}
\begin{itemize}
  \item Led the ``Reversed Instance'' development, enabling more cost-effective cloud resource offerings.
  \item Optimized database queries, reducing production database CPU utilization to 1/20 of its original level.
  \item Tech Stack: Java, Spring Boot, Redis, MySQL/TDSQL, Kafka, Docker, Kubernetes, Tencent Cloud
\end{itemize}

\job{WanGuoyou (Shenzhen) Technology Co.\ Ltd}{Software Engineer \arrow\ Software Manager}{Jul 2018 -- Jul 2021}
\begin{itemize}
  \item Joined as a backend software engineer on the founding team to build an online travel agency (OTA) platform.
  \item Promoted to lead the technology department within 1.5 years, leading the team of 10+ engineers.
  \item Streamlined core workflows from flight booking to financial operations. Improved operational efficiency by 10x.
  \item Integrated APIs from 10+ suppliers (flights, trains, hotels, car rentals), significantly expanding product offerings.
  \item Tech Stack: Java, Scala, Spring Boot, Spring Cloud, Redis, MySQL, RocketMQ, Kubernetes, Alibaba Cloud
\end{itemize}

% ---------- Honors & Awards ----------
\section{Honors \& Awards}
\textsc{Oregon Programming Languages Summer School Fellowship} | Jun 2024\par
\textsc{National Olympiad in Informatics in Provinces (NOIP), Third Prize, Sichuan, China} | \textsc{Oct 2010}\par

\end{document}
